% sec-05-proposed-clinical-research.tex - Proposed Clinical Research (final-ind)
\section{Proposed clinical research}

The proposed clinical research is presented in this section of the investigational
new drug (IND) application in the form required by 21 CFR \S 312.23(a)(6). It
consists of the study protocol for the planned Phase 1 first-in-human
investigation, the informed consent that governs each participant's enrollment,
and the investigator and facilities information that establishes that the study
will be conducted by qualified persons at adequate facilities. The investigation
described here is a prospective, open-label, single-arm, first-in-human study that
combines a standard 3+3 dose-finding evaluation of perioperative daraxonrasib
(RMC-6236) with an early-feasibility evaluation of an on-premises large language
model (LLM) directed eight-arm robotic pancreaticoduodenectomy (Whipple) system, a
significant-risk Class II collaborative device. Because the device participates in
the delivery of the investigational drug, the study is conducted under a combined
IND (21 CFR \S 312) and investigational device exemption (IDE, 21 CFR \S 812)
submission with a single integrated protocol, a single informed consent carrying
the Physical AI opt-out, and a single integrated safety governance structure. The
full clinical protocol, version 1.0.0, is bound to this IND as a supporting
document and is archived with the digital object identifier
DOI 10.5281/zenodo.20780121
(\href{https://doi.org/10.5281/zenodo.20780121}{https://doi.org/10.5281/zenodo.20780121});
the synopsis and supporting tables in \S 6.1 summarize that bound document, and
every quantitative value reported here is traceable to it
\cite{daraxwhipple,onpremwhippl}.

This section is written so that an FDA reviewer, an institutional review board (IRB)
member, a data safety monitoring board (DSMB) member, or an independent monitor can
read the proposed clinical research in isolation and reconstruct exactly how the
study is conducted, how each participant is protected, and how every regulatory
clock is met. The regulatory architecture is deliberately layered: the drug
component satisfies the content requirements of 21 CFR \S 312.23(a)(6) for a Phase 1
study protocol, namely a statement of objectives and purpose, the names and
addresses and qualifications of investigators, the criteria for participant
selection and exclusion and an estimate of the number to be studied, a description
of the design including the kind of control group, the method for determining doses
to be administered, a description of observations and measurements to fulfill the
objectives, and a description of the clinical procedures and laboratory tests
undertaken to monitor the effects of the drug and to minimize risk. The device
component layers on the IDE content of 21 CFR \S 812.25, namely the investigational
plan, the report of prior investigations, and the monitoring procedures, and the
Physical AI System Description meeting the eleven fields of 21 CFR \S 312.23(g) as
adapted for autonomy-bearing investigations. Throughout, the convention is that the
drug and the device are governed by one integrated protocol so that no safety signal
on either side is reviewed in isolation from the other, and that the autonomy-bearing
behavior of the system is held to a continuous human-oversight standard that prohibits
full autonomy at this phase consistent with 21 CFR \S 312.21(e)
\cite{phase1guidance,congress9510}.

\subsection{Study Protocol}

\subsubsection*{Protocol identification and overview}

The bound protocol is entitled the LLM-Directed PDAC Robotic Daraxonrasib Phase 1
study (protocol version 1.0.0, DOI 10.5281/zenodo.20780121). It is a prospective,
open-label, single-arm, first-in-human dose-finding and early-feasibility study in
adults with histologically or cytologically confirmed, resectable or
borderline-resectable, KRAS G12-mutated pancreatic ductal adenocarcinoma (PDAC)
who are candidates for pancreaticoduodenectomy. The central hypothesis is that the
combined drug-device intervention can be delivered with an acceptable early safety
profile and can complete the assigned operative tasks under continuous human
oversight while a recommended Phase 2 dose (RP2D) of daraxonrasib is identified,
supporting progression to a later-phase efficacy study. Enrollment is planned for
up to $n=18$ treated participants (approximately 36 screened) at a single academic
medical center, assigned across three daraxonrasib dose levels by the 3+3 rule
with staggered sentinel enrollment, at a planned accrual of approximately one to
two participants per month. The investigation is structured around the
early-feasibility pathway recognized for significant-risk devices, which generates
the device-performance and safety evidence required before any controlled
comparison is undertaken \cite{phase1guidance,daraxwhipple}.

The choice of a first-in-human dose-finding and early-feasibility design follows
directly from the regulatory status of the two components. Daraxonrasib has not
previously been studied in a perioperative schedule and has not previously been
paired with an autonomy-bearing surgical device, so no prior human experience
quantifies the safety of the combination; the device has never operated on a human
participant; and the integrated drug-device combination has been characterized only
in computational evidence. A small, intensively monitored, dose-escalating cohort
with staggered sentinel enrollment and binding halt rules is therefore the design
that minimizes the number of participants exposed before each successive safety
question is answered, consistent with the FDA guidance on first-in-human and Phase 1
oncology studies and with the early-feasibility framework for significant-risk
devices \cite{phase1guidance,oncotrialllm}. The disease context reinforces the
ethical case for a small, carefully escalated study: PDAC carries a total five-year
survival below 13\% and is the third leading cause of cancer death in the United
States and the fourth in Europe, so the curative-intent operation under study is
performed in a population for whom every avoidable surgical delay is consequential,
and for whom the alternative to enrollment is the same demanding operation without
the investigational advisory layer \cite{pdac060s2030}.

\subsubsection*{Objectives and endpoints}

The protocol specifies three co-primary objectives, each evaluated with exact
two-sided 95\% Clopper-Pearson confidence intervals because the design is
descriptive and non-confirmatory under ICH E9 with no formal power calculation and
no multiplicity adjustment. The first co-primary endpoint is the incidence of
device- or procedure-related serious adverse events through 30 days, including any
Clavien-Dindo Grade III or higher complication. The second co-primary endpoint is
the maximum tolerated dose (MTD) and the RP2D of perioperative daraxonrasib,
determined by the 3+3 rule across a 28-day dose-limiting toxicity (DLT) window. The
third co-primary endpoint is technical feasibility, defined as the proportion of
participants completing the assigned operative tasks without unplanned conversion
attributable to device malfunction or unsafe device behavior. The secondary
endpoints are the margin-negative (R0) resection rate on day-90 pathology
(per-protocol), the International Study Group of Pancreatic Surgery (ISGPS) Grade B
or C postoperative pancreatic-fistula rate, the Clavien-Dindo Grade III or higher
complication rate at 90 days, 30-day and 90-day mortality, the conversion rate, and
operative parameters (estimated blood loss, operating-room time, and length of
stay), each benchmarked descriptively against the 2025 Dutch nationwide robotic
pancreaticoduodenectomy cohort (conversion 10.1\%, Clavien-Dindo III+ 41.3\%, ISGPS
B/C 24.4\%, 30-day mortality 3.9\%). Exploratory endpoints are progression-free and
overall survival to 24 months by RECIST version 1.1 with Kaplan-Meier estimation,
the concordance of the LLM advisory with the hard-coded safety gate, and the
sim-to-real gap in millimeters and newtons \cite{daraxwhipple,oncotrialllm}.

Each endpoint is paired with an explicit definition, an analysis population, and a
prespecified summary statistic so that the analysis is fully prespecified before the
first participant is enrolled. Table~\ref{tab:sec06-endpoints} sets out the complete
endpoint hierarchy with the definition, the population, and the summary measure for
each endpoint. The three co-primary endpoints are deliberately heterogeneous: a
device-and-procedure safety endpoint, a drug dose-finding endpoint, and a device
feasibility endpoint. This structure reflects the combined IND/IDE nature of the
investigation, in which the dose-finding question for the drug and the readiness and
feasibility questions for the device must be answered together, because the device
participates in delivering the drug and the perioperative restart of the drug depends
on the operative and pathology results of the device procedure.

\needspace{10\baselineskip}
\begin{table}[t]
\caption{Table 6.1. Endpoint hierarchy for the bound Phase 1 protocol, with the
definition, the analysis population, and the prespecified summary measure for each
co-primary, secondary, and exploratory endpoint. All binomial endpoints are
summarized with exact two-sided 95\% Clopper-Pearson confidence intervals.}
\label{tab:sec06-endpoints}
\centering
\begin{tabularx}{\textwidth}{L{3.0cm} Y L{2.6cm} L{2.6cm}}
\toprule
\textbf{Endpoint} & \textbf{Definition} & \textbf{Population} & \textbf{Summary measure} \\
\midrule
Co-primary 1: device/procedure SAE & Incidence of device- or procedure-related serious adverse events through 30 days, including any Clavien-Dindo Grade III+ complication & Safety (all dosed) & Proportion with exact 95\% CI \\
Co-primary 2: MTD and RP2D & Maximum tolerated dose and recommended Phase 2 dose by the 3+3 rule across the 28-day DLT window & DLT-evaluable & Dose level; DLT rate per level with exact 95\% CI \\
Co-primary 3: technical feasibility & Proportion completing assigned operative tasks without unplanned conversion from device malfunction or unsafe behavior & Safety (all dosed) & Proportion with exact 95\% CI \\
Secondary: R0 resection & Margin-negative resection on day-90 synoptic pathology & Per-protocol & Proportion with exact 95\% CI \\
Secondary: ISGPS B/C fistula & Grade B or C postoperative pancreatic fistula by the ISGPS definition & Safety (all dosed) & Proportion vs Dutch 24.4\% \\
Secondary: Clavien-Dindo III+ & Grade III or higher complication at 90 days & Safety (all dosed) & Proportion vs Dutch 41.3\% \\
Secondary: mortality & 30-day and 90-day mortality & Safety (all dosed) & Proportion vs Dutch 3.9\% (30-day) \\
Secondary: conversion & Conversion to open or manual technique & Safety (all dosed) & Proportion vs Dutch 10.1\% \\
Secondary: operative parameters & Estimated blood loss, operating-room time, length of stay & Safety (all dosed) & Median with range \\
Exploratory: PFS and OS & Progression-free and overall survival to 24 months by RECIST v1.1 & mITT & Kaplan-Meier estimate \\
Exploratory: advisory concordance & Concordance of the LLM advisory with the hard-coded safety gate & Safety (all dosed) & Proportion of concordant verdicts \\
Exploratory: sim-to-real gap & Trajectory gap (mm) and tip-force gap (N) between simulation and the operative record & Safety (all dosed) & Median (mm; N) \\
\bottomrule
\end{tabularx}
\end{table}

Table~\ref{tab:sec06-soa} presents the schedule of activities that governs the
timing of every assessment, from screening through the operative episode and the
24-month overall-survival follow-up.

\needspace{10\baselineskip}
\begin{table}[t]
\caption{Table 6.2. Schedule of activities (summary) for the bound Phase 1 protocol.
The intensive safety-collection window spans consent through the 30-day
postoperative visit and encompasses the 28-day daraxonrasib dose-limiting toxicity
window.}
\label{tab:sec06-soa}
\centering
\begin{tabularx}{\textwidth}{L{4.3cm} C{1.4cm} C{1.6cm} C{1.6cm} Y}
\toprule
\textbf{Assessment} & \textbf{Screen} & \textbf{Pre-op / day 0} & \textbf{Post-op to day 30} & \textbf{Long-term follow-up} \\
\midrule
Informed consent and Physical AI opt-out election & X & - & - & re-consent on amendment \\
Eligibility, ECOG, KRAS G12 genotype, staging imaging & X & - & - & - \\
Multidisciplinary resectability review & X & X & - & - \\
Daraxonrasib pre-operative course (oral once daily) & - & X & - & - \\
Protocol-mandated operative pause of daraxonrasib & - & X & - & - \\
Robotic Whipple (8 operative phases P1 to P8) & - & X & - & - \\
Pre-procedure safety matrix and USL confirmation & - & X & - & re-check after any change \\
Perioperative pause-and-restart advisory (T+7/14/21d) & - & - & X & - \\
Vitals, hematology, chemistry, hepatic/pancreatic indices & X & X & X & per follow-up schedule \\
Serum daraxonrasib trough assay & - & X & X & as indicated \\
ISGPS fistula grading (PJ dominant) & - & - & X & - \\
Adverse-event and Physical AI AE capture & X & X & X & SAE/PAI AE to 24 months \\
RECIST v1.1 imaging (central read) & X & - & X & per follow-up schedule \\
Pathology, R0 margin (day-90 read, per-protocol) & - & - & X & day 90 \\
Progression-free and overall survival & - & - & - & to 24 months (q12 weeks) \\
\bottomrule
\end{tabularx}
\end{table}

\subsubsection*{Study design and dose escalation}

The drug component proceeds under the IND pathway and the device component under
the IDE pathway, with a single integrated protocol and a single safety governance
structure. The escalation uses three dose levels anchored to the pan-KRAS RASolute
302 program, in which 300 mg once daily is the established RP2D: DL1 is 160 mg once
daily (starting dose, chosen below the reference dose to preserve a conservative
first-in-human margin), DL2 is 220 mg once daily, and DL3 is the 300 mg once-daily
reference dose. Three participants are enrolled at each level and observed across a
28-day DLT window. If no DLT occurs among the first three, the cohort escalates; if
one of three experiences a DLT, the cohort expands to six; a single DLT among six
permits escalation, whereas two or more DLTs at a level fix the MTD at the
preceding level. The MTD is the highest level at which fewer than two of six
DLT-evaluable participants experience a DLT, and the RP2D is selected at or below
the MTD. With a maximum of three levels of six, the design treats up to 18
participants. The analysis populations are the Safety population (all dosed
participants), the DLT-evaluable population, the Per-protocol population, and the
modified intention-to-treat (mITT) population.

The 3+3 decision rule is applied uniformly at every level, and the precise transition
logic is set out in Table~\ref{tab:sec06-33}. The rule is intentionally simple and
fully deterministic so that the escalation decision can be reproduced by any reviewer
from the observed DLT count alone, and so that no judgment beyond the DLT
adjudication is required to choose the next cohort. A DLT is any drug-attributable
toxicity meeting the protocol grading thresholds during the 28-day window after the
first dose at a given level; events adjudicated as attributable to the surgery or the
device rather than to the drug are captured under the device/procedure safety endpoint
and the Physical AI safety stream rather than as DLTs, so that the dose-finding
question for the drug is not contaminated by device-attributable events.

\needspace{9\baselineskip}
\begin{table}[t]
\caption{Table 6.3. The 3+3 dose-escalation decision rule applied at each
daraxonrasib dose level (DL1 160 mg, DL2 220 mg, DL3 300 mg once daily) across the
28-day dose-limiting toxicity window. The maximum tolerated dose is the highest
level at which fewer than two of six dose-limiting-toxicity-evaluable participants
experience a dose-limiting toxicity.}
\label{tab:sec06-33}
\centering
\begin{tabularx}{\textwidth}{L{3.4cm} L{3.4cm} Y}
\toprule
\textbf{DLTs in first 3} & \textbf{Action} & \textbf{Resolution} \\
\midrule
0 of 3 & Escalate to the next dose level & Current level cleared; move up one level (or declare RP2D at DL3) \\
1 of 3 & Expand the current level to 6 & Enroll 3 more at the same level and re-evaluate at 6 \\
2 or 3 of 3 & Stop escalation at this level & MTD exceeded; MTD is the preceding level \\
1 of 6 (after expansion) & Escalate to the next dose level & Level tolerated; move up one level (or declare RP2D at DL3) \\
2 or more of 6 & Stop escalation at this level & MTD exceeded; MTD is the preceding level \\
\bottomrule
\end{tabularx}
\end{table}

The estimation precision that this sample supports is illustrated by the exact
Clopper-Pearson intervals, which are the appropriate interval for a small-sample
binomial proportion because they guarantee at least nominal coverage without relying
on a normal approximation. An observed rate of 0 of 18 yields 0\% (95\% CI 0 to
18\%), 1 of 18 yields 5.6\% (95\% CI 0.1 to 27\%), 0 of 3 yields 0\% (95\% CI 0 to
71\%), and 0 of 6 yields 0\% (95\% CI 0 to 46\%) \cite{daraxwhipple}. These worked
values are reproduced in Table~\ref{tab:sec06-cp} so that a reviewer can read off the
inferential reach of each cohort size directly. The contrast between the 0-of-3 and
0-of-18 rows makes the rationale for staggered sentinel enrollment explicit: a single
cleared cohort of three bounds the true serious-event rate only loosely (up to 71\%),
whereas the full 18-participant experience tightens the upper bound to 18\%, so the
study deliberately defers any claim of acceptable safety until the cumulative
experience supports it.

\needspace{8\baselineskip}
\begin{table}[t]
\caption{Table 6.4. Worked exact two-sided 95\% Clopper-Pearson confidence intervals
for representative observed event counts at the cohort and full-study sample sizes.
The intervals quantify the descriptive precision of the non-confirmatory design and
motivate the staggered sentinel enrollment.}
\label{tab:sec06-cp}
\centering
\begin{tabularx}{\textwidth}{C{2.6cm} C{2.6cm} C{2.6cm} Y}
\toprule
\textbf{Events} & \textbf{Denominator} & \textbf{Point estimate} & \textbf{Exact two-sided 95\% Clopper-Pearson interval} \\
\midrule
0 & 3 & 0\% & 0 to 71\% \\
0 & 6 & 0\% & 0 to 46\% \\
0 & 18 & 0\% & 0 to 18\% \\
1 & 18 & 5.6\% & 0.1 to 27\% \\
\bottomrule
\end{tabularx}
\end{table}

The device has no pharmacologic dose; its dose analogue is the readiness evidence
required before any patient-facing activity, escalated only by passing the
mandatory Phase 0 simulation-validation gate. Three conditions must be met and
documented before the first robotic case: a Unified Safety Level (USL) rating of at
least 7.0 on the surgical-readiness scale; at least 1000 simulated procedures
completed across at least two independent simulation frameworks (for example Isaac,
MuJoCo, Gazebo, or PyBullet) so that the readiness finding is not an artifact of any
single simulator; and a sim-to-real fidelity within prespecified limits, a
trajectory gap of less than 2 mm and a tip-force gap of less than 0.5 N between the
simulation prediction and the operative record. The gate is re-run after any
software, model, or configuration change before any further patient-facing use, so
that no model update, retraining, or reconfiguration reaches a participant without a
fresh demonstration of readiness \cite{onpremwhippl,fdadigtwinpc}.
Table~\ref{tab:sec06-phase0} states the three Phase 0 gate conditions, their numeric
thresholds, and the consequence of a failed condition.

\needspace{8\baselineskip}
\begin{table}[t]
\caption{Table 6.5. The mandatory Phase 0 simulation-validation gate conditions that
constitute the device readiness analogue of a drug dose. All three conditions must
be met and documented before the first robotic case, and the gate is re-run after
any software, model, or configuration change.}
\label{tab:sec06-phase0}
\centering
\begin{tabularx}{\textwidth}{L{4.0cm} L{4.6cm} Y}
\toprule
\textbf{Condition} & \textbf{Threshold} & \textbf{Consequence if not met} \\
\midrule
Unified Safety Level (USL) rating & At least 7.0 on the surgical-readiness scale & No patient-facing robotic activity; remediation and re-rating required \\
Simulated procedure volume & At least 1000 simulated procedures across at least 2 independent frameworks (Isaac, MuJoCo, Gazebo, PyBullet) & Readiness finding not established; additional simulation required before any case \\
Sim-to-real fidelity & Trajectory gap less than 2 mm and tip-force gap less than 0.5 N versus the operative record & Gate fails; the system may not proceed and is re-validated after correction \\
\bottomrule
\end{tabularx}
\end{table}

\subsubsection*{Population and eligibility}

The study population comprises adults with histologically or cytologically confirmed,
resectable or borderline-resectable, KRAS G12-mutated PDAC who are eligible for
perioperative daraxonrasib under the broad pan-KRAS RASolute 302 criteria and whose
tumor and vascular anatomy are compatible with the eight-arm robotic system.
Participant selection is equitable; eligibility is determined solely by the
scientific and safety criteria, and limited English proficiency is not an exclusion
criterion. The synthetic reference exemplar PAT-PDAC-0001 (a 68-year-old with a
3.4 cm head-of-pancreas tumor abutting the superior mesenteric vein at $75^\circ$,
KRAS G12D, microsatellite stable, CA 19-9 of 412 U/mL, ECOG performance status 1,
after a completed neoadjuvant modified FOLFIRINOX course of four cycles) illustrates
a representative eligible participant and does not represent an enrolled individual.
The full eligibility criteria are summarized in Table~\ref{tab:sec06-eligibility}.

\needspace{10\baselineskip}
\begin{table}[t]
\caption{Table 6.6. Key inclusion and exclusion criteria for the bound Phase 1
protocol.}
\label{tab:sec06-eligibility}
\centering
\begin{tabularx}{\textwidth}{Y Y}
\toprule
\textbf{Inclusion criteria (all required)} & \textbf{Exclusion criteria (any one excludes)} \\
\midrule
Signed informed consent before any study-specific procedure, including explicit Physical AI opt-out election or declination & Distant metastatic disease on staging \\
Age 18 years or older at consent & Vascular involvement precluding a margin-negative (R0) resection, or anatomy incompatible with the eight-arm robotic approach \\
Histologically or cytologically confirmed resectable or borderline-resectable PDAC on multidisciplinary review & Prior pancreatic resection or prior surgery precluding the planned reconstruction \\
Documented KRAS G12 mutation (for example G12D, G12V, or G12R) on a validated assay & Uncontrolled intercurrent illness or comorbidity making major surgery or daraxonrasib unsafe \\
ECOG performance status 0 or 1 & Pregnancy or lactation, or unwillingness to use adequate contraception where applicable \\
Adequate organ and marrow function for major pancreaticobiliary surgery and daraxonrasib & Known contraindication or hypersensitivity to daraxonrasib or a required excipient, or an unmanageable drug interaction \\
Measured baseline CA 19-9 recorded for serial assessment; eligible for daraxonrasib under RASolute 302 pan-KRAS criteria & Inability to provide informed consent or to comply with the Schedule of Activities \\
Tumor and vascular anatomy compatible with the eight-arm robotic system on preoperative imaging & Anatomy or comorbidity otherwise incompatible with safe conduct of the combined intervention \\
\bottomrule
\end{tabularx}
\end{table}

The eligibility criteria are constructed so that each criterion maps to an identified
risk. The KRAS G12 genotype requirement establishes the molecular target for
daraxonrasib and aligns the population with the pan-KRAS RASolute 302 reference
program from which the dose levels are drawn; the resectable or
borderline-resectable staging and the multidisciplinary resectability review ensure
that the curative-intent operation is appropriate and that vascular involvement
precluding an R0 resection is excluded before any device exposure; the ECOG 0 or 1
and adequate organ and marrow function requirements ensure that the participant can
tolerate both the major pancreaticobiliary operation and the RAS(ON) inhibitor; and
the anatomic-compatibility requirement ensures that the tumor and the five named
vessels are within the operating envelope of the eight-arm platform. Limited English
proficiency is explicitly not an exclusion criterion, and qualified interpreters and
translated materials are provided, so that the equitable-selection principle of the
governing human-subjects regulations is met and the population is not narrowed in a
way that would compromise generalizability or fairness \cite{cfr050paiuni}.

\subsubsection*{Intervention: drug and device}

Daraxonrasib is an orally bioavailable RAS(ON) multi-selective (pan-RAS)
small-molecule inhibitor that binds the active, guanosine-triphosphate-bound state
of mutant RAS in complex with cyclophilin A, suppressing downstream effector
signaling across the common KRAS variants in PDAC; it received FDA Breakthrough
Therapy designation in June 2025 for previously treated metastatic KRAS G12-mutant
PDAC. In this study it is administered orally once daily in a perioperative
schedule comprising a defined pre-operative course, a protocol-mandated pause across
the operative window, and a postoperative restart governed by the advisory in
\S 6.1 below. The device component is an eight-arm parallel coelomic surgical
platform carrying 56 mechanical degrees of freedom (8 arms $\times$ 7 per-arm
degrees of freedom) and a 640-channel mixed-rate sensor stack (80 channels per arm),
with force channels sampling at 100 kHz and all remaining channels at the 10 kHz
command rate, and a positioning accuracy of 0.05 mm root mean square at a peak tip
velocity of 1,200 mm/s. An on-premises repository-pinned, open-weights LLM acts as a
second-opinion advisory oracle held strictly outside the robot vendor kinematic
stack and hash-pinned to a deterministic seed and a repository commit; it is never
an autonomous controller, and full autonomy is prohibited at this phase consistent
with 21 CFR \S 312.21(e). The advisory is gated by a verification-before-generation
ten-gate suite returning ACCEPT, BLOCK, or ESCALATE, aligned with H.R. 9510, the
Verification Before Generation in Physical AI Oncology Trials Act of 2026
\cite{congress9510,onpremwhippl}.

The hard cyber-physical safety envelope comprises a per-arm tip-force cap of
$\leq 3$ N, a cumulative cross-arm tip-force cap of $\leq 18$ N enforced at the
heartbeat boundary, a 10 kHz heartbeat coordination bus broadcasting a 64-byte frame
on a $100\,\mu$s per-arm watchdog deadline, an emergency arm park within $50\,\mu$s
of a missed or out-of-parameter frame, a cross-arm emergency stop (E-stop) that
halts the full platform in $\leq 3$ ms at a 0.0 N force clamp, and a
software-independent hardware E-stop meeting the 21 CFR \S 312.404 requirement of
$\leq 500$ ms. A vascular no-fly gate sampled at 10 kHz evaluates instrument-tip
proximity to five named vessels (superior mesenteric vein, portal vein, hepatic
artery, celiac axis, and superior mesenteric artery) against three concentric
thresholds: a soft-warning zone at 3.0 mm raises an LLM advisory, a no-fly zone
(1.0 mm at the superior mesenteric and portal veins; 1.5 mm at the hepatic artery,
celiac axis, and superior mesenteric artery) blocks the command, and a hard-stop
zone at 5.0 mm forces the $\leq 3$ ms E-stop; across a representative 100-tick
verdict sample the gate returned clear on 83 ticks, soft-warning on 6, no-fly on 6,
and hard-stop on 5. The operation proceeds across eight operative phases: phases P1
through P4 comprise dissection and specimen removal, phase P5 the duct-to-mucosa
pancreaticojejunostomy under a controller ring-tension band of 0.30 to 0.60 N, phase
P6 the end-to-side hepaticojejunostomy at 0.20 to 0.50 N, phase P7 the antecolic
gastrojejunostomy at 0.40 to 0.80 N, and phase P8 final imaging, hemostasis, and
drain placement \cite{onpremwhippl}.

The numeric parameters of the safety envelope are consolidated in
Table~\ref{tab:sec06-envelope} so that a reviewer can confirm that every quoted
limit, latency, and threshold is internally consistent and traceable to the bound
protocol. The envelope is structured as a defense-in-depth chain in which a soft
advisory layer (the 3.0 mm soft-warning zone and the LLM second opinion) is backed
by a command-blocking layer (the no-fly zones and the force caps) which is in turn
backed by an irreversible halt layer (the 5.0 mm hard-stop zone, the $\leq 3$ ms
cross-arm E-stop, and the software-independent hardware E-stop). No single layer is
relied upon alone, and the hardware E-stop is independent of all software so that a
software fault cannot disable the final halt path. The per-phase controller
ring-tension bands constrain the anastomotic forces during the reconstruction phases
P5 through P7, where the surgical risk of a leak is concentrated, and the resulting
pancreaticojejunostomy fistula grade is the dominant input to the perioperative
restart advisory and the drug stopping rule.

\needspace{10\baselineskip}
\begin{table}[t]
\caption{Table 6.7. The hard cyber-physical safety envelope of the eight-arm robotic
platform, with the numeric value and the regulatory or operative role of each
parameter. Every value is traceable to the bound protocol and to the Physical AI
System Description.}
\label{tab:sec06-envelope}
\centering
\begin{tabularx}{\textwidth}{L{4.6cm} L{3.4cm} Y}
\toprule
\textbf{Parameter} & \textbf{Value} & \textbf{Role} \\
\midrule
Mechanical degrees of freedom & 56 (8 arms $\times$ 7) & Operating envelope of the platform \\
Sensor channels & 640 (80 per arm) & 100 kHz force; 10 kHz all other \\
Positioning accuracy & 0.05 mm RMS at 1,200 mm/s & Trajectory specification benchmarked by the sim-to-real gap \\
Per-arm tip-force cap & $\leq 3$ N & Hard force limit per arm \\
Cumulative cross-arm cap & $\leq 18$ N & Hard force limit across arms, enforced at the heartbeat boundary \\
Heartbeat coordination bus & 10 kHz, 64-byte frame & $100\,\mu$s per-arm watchdog deadline \\
Emergency arm park & within $50\,\mu$s & On a missed or out-of-parameter frame \\
Cross-arm E-stop & $\leq 3$ ms at 0.0 N clamp & Halts the full platform \\
Hardware E-stop & $\leq 500$ ms & Software-independent; 21 CFR \S 312.404 \\
Soft-warning vascular zone & 3.0 mm & Raises an LLM advisory \\
No-fly vascular zone & 1.0 mm (SMV, PV); 1.5 mm (HA, CA, SMA) & Blocks the command \\
Hard-stop vascular zone & 5.0 mm & Forces the $\leq 3$ ms E-stop \\
\bottomrule
\end{tabularx}
\end{table}

\subsubsection*{Perioperative daraxonrasib pause-and-restart advisory}

The two components are coupled at one decision point only: the perioperative
pause-and-restart advisory for the drug is informed by the operative and pathology
results of the device procedure. Daraxonrasib is paused before surgery and restarted
postoperatively on a day chosen by an advisory that is LLM-bound and human-confirmed,
never autonomous. The advisory takes two inputs, the pancreaticojejunostomy ISGPS
fistula grade (A, B, or C) and the serum daraxonrasib trough relative to a 0.5 ng/mL
threshold, and recommends a restart at postoperative day T+7d, T+14d, or T+21d (or a
continued hold). In the 32-iteration deterministic simulation sweep, baseline serum
was 0.45 ng/mL at the operative timepoint across all iterations (below the 0.5 ng/mL
trough threshold), and the advisory recommended T+7d in 29 of 32 iterations (90.6\%,
Grade A fistula with sub-threshold trough), T+14d in 3 of 32 (9.4\%, Grade B or
delayed serum rise), and T+21d or continued hold in 0 of 32 (reserved for Grade C,
which also maps to the drug stopping rule). Figure~15 reproduces the advisory logic
and the sweep distribution.

The advisory illustrates the disciplined separation between the autonomy-bearing
component and the human decision: the LLM proposes a restart day from the two inputs,
but the restart is executed only after a qualified clinician confirms it, so that no
drug-administration decision is ever made autonomously. The mapping is monotone in
both inputs, with a lower fistula grade and a sub-threshold trough favoring the
earliest restart (T+7d) and a higher grade or a rising trough favoring a later
restart (T+14d) or a hold; a confirmed Grade C fistula precludes restart entirely and
maps to the drug stopping rule. Table~\ref{tab:sec06-advisory} sets out the worked
input-to-output mapping together with the sweep frequencies, so that the advisory can
be audited against the deterministic sweep distribution reproduced in the figure.

\needspace{9\baselineskip}
\begin{table}[t]
\caption{Table 6.8. The perioperative daraxonrasib pause-and-restart advisory mapping
and the 32-iteration deterministic sweep distribution. The two inputs are the
pancreaticojejunostomy ISGPS fistula grade and the serum daraxonrasib trough relative
to the 0.5 ng/mL threshold; the output is a human-confirmed restart day. Baseline
serum was 0.45 ng/mL (sub-threshold) across all iterations.}
\label{tab:sec06-advisory}
\centering
\begin{tabularx}{\textwidth}{L{3.0cm} L{4.2cm} L{2.8cm} Y}
\toprule
\textbf{Restart output} & \textbf{Input condition} & \textbf{Sweep frequency} & \textbf{Disposition} \\
\midrule
T+7d & Grade A fistula with sub-threshold trough (below 0.5 ng/mL) & 29 of 32 (90.6\%) & Earliest restart; resume daraxonrasib \\
T+14d & Grade B fistula or a delayed serum rise & 3 of 32 (9.4\%) & Deferred restart; resume after re-assessment \\
T+21d or continued hold & Grade C fistula & 0 of 32 & Restart precluded; maps to the drug stopping rule \\
\bottomrule
\end{tabularx}
\end{table}

\needspace{8\baselineskip}
\begin{mermaidfig}
\node[mmin]  (pause) at (0,0)        {Daraxonrasib paused\\before surgery};
\node[mmin]  (in1)   at (0,-2.0)     {Input 1: PJ ISGPS\\fistula grade A / B / C};
\node[mmin]  (in2)   at (0,-4.0)     {Input 2: serum trough\\vs 0.5 ng/mL threshold};
\node[mmproc] (llm)  at (5.6,-2.0)   {LLM advisory\\(hash-pinned)\\+ human confirmation\\(never autonomous)};
\node[mmgoal] (t7)   at (11.2,0)     {Restart T+7d\\29 of 32 (Grade A,\\sub-threshold trough)};
\node[mmstep] (t14)  at (11.2,-2.0)  {Restart T+14d\\3 of 32 (Grade B or\\delayed serum rise)};
\node[mmdark] (t21)  at (11.2,-4.0)  {Restart T+21d or hold\\0 of 32 (Grade C)};
\draw[mmedge] (pause) to[out=0,in=180,looseness=0.6] (llm.west);
\draw[mmedge] (in1) -- (llm.west);
\draw[mmedge] (in2) to[out=0,in=180,looseness=0.6] (llm.west);
\draw[mmedge] (llm) to[out=0,in=180,looseness=0.6] (t7.west);
\draw[mmedge] (llm) -- (t14.west);
\draw[mmedge] (llm) to[out=0,in=180,looseness=0.6] (t21.west);
\end{mermaidfig}
\figcaption{Figure 15. The perioperative daraxonrasib pause-and-restart advisory.
The drug is paused before surgery and the postoperative restart day is set by a
hash-pinned, human-confirmed LLM advisory (never autonomous) from two inputs: the
pancreaticojejunostomy ISGPS fistula grade (A, B, or C) and the serum daraxonrasib
trough relative to a 0.5 ng/mL threshold. In the 32-iteration deterministic sweep
the advisory recommended restart at T+7 days in 29 of 32 iterations (Grade A with
sub-threshold trough), T+14 days in 3 of 32 (Grade B or delayed serum rise), and
T+21 days or continued hold in 0 of 32 (reserved for Grade C).}

\subsubsection*{Assessments and the vascular and halt safety chain}

Clinical safety assessments comprise vital signs, a directed physical examination,
hematology, serum chemistry including hepatic and pancreatic indices, coagulation
studies, the serum daraxonrasib trough assay, and protocol-specified imaging, at the
frequencies in the schedule of activities. In addition, the device generates a
continuous safety telemetry stream that is a source of protocol observations under
21 CFR \S 312.23(g): per-arm and cumulative force profiles, trajectory accuracy
against the 0.05 mm specification, every vascular safety-zone verdict, every
heartbeat and watchdog event, and every E-stop activation, each captured at its
recorded rate in the hash-chained audit trail. The halt chain is layered: the 10 kHz
heartbeat bus broadcasts a 64-byte frame on a $100\,\mu$s per-arm watchdog deadline;
a frame that is on time continues the command state, whereas a missed or
out-of-parameter frame parks the affected arm within $50\,\mu$s and escalates to the
cross-arm E-stop that halts the platform in $\leq 3$ ms at a 0.0 N force clamp,
backed by the software-independent hardware E-stop. Postoperative pancreatic fistula
is graded by the ISGPS definition as a biochemical leak, Grade B (requiring a change
in management), or Grade C (requiring reoperation or causing organ failure or death),
and the pancreaticojejunostomy grade is the dominant determinant and the input to the
restart advisory and the drug stopping rule \cite{daraxwhipple,oncotrialllm}.

The vascular no-fly gate is the most safety-critical of the telemetry-driven
controls, because the five named vessels border the operative field of a
pancreaticoduodenectomy and an inadvertent vascular injury is among the most lethal
intraoperative events. The gate is evaluated at 10 kHz against the three concentric
thresholds, and the representative 100-tick verdict sample (clear 83, soft-warning 6,
no-fly 6, hard-stop 5) demonstrates the expected distribution in which most ticks are
clear, a minority raise a soft advisory, a smaller minority block a command, and the
hard-stop verdicts force the E-stop. Table~\ref{tab:sec06-vascular} sets out the three
vascular zones, their per-vessel thresholds, the system action, and the corresponding
verdict count in the representative sample, so that the proximity controls are fully
specified and auditable against the telemetry record.

\needspace{8\baselineskip}
\begin{table}[t]
\caption{Table 6.9. The vascular no-fly gate sampled at 10 kHz against five named
vessels, with the per-vessel proximity threshold, the system action, and the verdict
count in a representative 100-tick sample. SMV, superior mesenteric vein; PV, portal
vein; HA, hepatic artery; CA, celiac axis; SMA, superior mesenteric artery.}
\label{tab:sec06-vascular}
\centering
\begin{tabularx}{\textwidth}{L{3.0cm} L{4.2cm} Y C{2.0cm}}
\toprule
\textbf{Zone} & \textbf{Threshold} & \textbf{System action} & \textbf{Sample verdicts} \\
\midrule
Clear & Beyond all thresholds & Command proceeds normally & 83 of 100 \\
Soft-warning & 3.0 mm (all five vessels) & Raises an LLM advisory; command continues under oversight & 6 of 100 \\
No-fly & 1.0 mm (SMV, PV); 1.5 mm (HA, CA, SMA) & Blocks the command & 6 of 100 \\
Hard-stop & 5.0 mm & Forces the $\leq 3$ ms cross-arm E-stop at 0.0 N & 5 of 100 \\
\bottomrule
\end{tabularx}
\end{table}

\subsubsection*{Safety reporting clocks and the six Physical AI triggers}

Adverse events are captured through two parallel and equally binding streams, a
conventional pharmacovigilance stream for clinical events graded by CTCAE version 5.0,
and a Physical AI safety stream for device- and AI-related events under 21 CFR
\S 312.32(f). Serious adverse events are reported by the investigator to the sponsor
within 24 hours of awareness, and the sponsor submits IND/IDE safety reports on the
regulatory clocks: any unexpected fatal or life-threatening suspected adverse
reaction is reported as soon as possible but no later than 7 calendar days after
initial receipt of the information, and any other serious and unexpected suspected
adverse reaction is reported no later than 15 calendar days after the determination
that the event is reportable (21 CFR \S 312.32). In addition, six Physical AI
reporting triggers are reported under \S 312.32(g): a serious Physical AI adverse
event on the 7- or 15-day timeline as applicable; a system-drug interaction causing
incorrect drug administration, regardless of seriousness, within 15 days; a
cybersecurity incident within 15 days (or 7 days if it has the potential to cause
patient harm); sustained AI model degradation below the pre-specified acceptance
criteria for three or more consecutive procedures or a cumulative 24 hours; a
sim-to-real divergence beyond twice the validated gap tolerance (trajectory deviation
greater than 4 mm or force discrepancy greater than 1.0 N); and a digital-twin
discrepancy beyond the pre-specified accuracy thresholds. For every Physical AI
adverse event, the complete hash-chained audit trail from 24 hours before the event
to 72 hours after the event is preserved under 21 CFR part 11 and made available to
the FDA on request. Figure~16 reproduces the clocks and the six triggers
\cite{phase1guidance,cfr312paiuni}.

The six Physical AI triggers extend conventional pharmacovigilance to the
autonomy-bearing failure modes that a drug-only IND does not contemplate, and each is
mapped to a defined clock and a defined evidentiary preservation in
Table~\ref{tab:sec06-triggers}. Two features distinguish the Physical AI stream from
the conventional stream. First, several triggers fire regardless of clinical
seriousness: a system-drug interaction causing incorrect drug administration is
reportable even if no harm results, and a cybersecurity incident is reportable on a
shortened 7-day clock when patient harm is merely possible, because the
autonomy-bearing context makes a near-miss a leading indicator of a future serious
event. Second, every Physical AI adverse event triggers preservation of the
hash-chained audit trail from 24 hours before to 72 hours after the event under 21
CFR part 11, so that the complete command, telemetry, advisory, and override record
surrounding the event is available for reconstruction. The two streams are equally
binding, and a single event that is both a clinical serious adverse event and a
Physical AI adverse event is reported on whichever clock is shorter.

\needspace{10\baselineskip}
\begin{table}[t]
\caption{Table 6.10. The two safety-reporting clocks and the six Physical AI
reporting triggers of 21 CFR \S 312.32(g), with the reporting clock and the
evidentiary preservation for each. Every Physical AI adverse event preserves the
hash-chained audit trail from minus 24 to plus 72 hours under 21 CFR part 11.}
\label{tab:sec06-triggers}
\centering
\begin{tabularx}{\textwidth}{L{4.8cm} L{3.2cm} Y}
\toprule
\textbf{Event} & \textbf{Reporting clock} & \textbf{Evidentiary preservation} \\
\midrule
Fatal or life-threatening suspected adverse reaction & 7 calendar days & Standard pharmacovigilance record \\
Other serious and unexpected suspected adverse reaction & 15 calendar days & Standard pharmacovigilance record \\
1. Serious Physical AI adverse event & 7 or 15 calendar days as applicable & Audit trail $-24$h to $+72$h \\
2. System-drug interaction causing incorrect drug administration & 15 calendar days (regardless of seriousness) & Audit trail $-24$h to $+72$h \\
3. Cybersecurity incident & 15 calendar days; 7 days if patient harm is possible & Audit trail $-24$h to $+72$h \\
4. Sustained model degradation & 15 calendar days; trigger at $\geq 3$ consecutive procedures or 24 cumulative hours & Audit trail $-24$h to $+72$h \\
5. Sim-to-real divergence beyond $2\times$ tolerance & 15 calendar days; trajectory over 4 mm or force over 1.0 N & Audit trail $-24$h to $+72$h \\
6. Digital-twin discrepancy beyond accuracy thresholds & 15 calendar days & Audit trail $-24$h to $+72$h \\
\bottomrule
\end{tabularx}
\end{table}

\needspace{8\baselineskip}
\begin{mermaidfig}
\node[mmin]  (ev)  at (0,-2.4)      {Adverse event\\(clinical or\\Physical AI)};
\node[mmgoal] (c7) at (5.4,0)       {7 calendar days\\fatal / life-threatening\\suspected adverse reaction};
\node[mmgoal] (c15) at (5.4,-2.4)   {15 calendar days\\other serious + unexpected\\suspected adverse reaction};
\node[mmproc] (g1) at (5.4,-4.4)    {1 Serious Physical AI AE (7 / 15d)};
\node[mmproc] (g2) at (5.4,-5.8)    {2 System-drug interaction (15d)};
\node[mmproc] (g3) at (10.8,-4.4)   {3 Cybersecurity incident (15d; 7d if harm)};
\node[mmproc] (g4) at (10.8,-5.8)   {4 Model degradation ($\geq 3$ procedures / 24h)};
\node[mmproc] (g5) at (5.4,-7.2)    {5 Sim-to-real divergence $> 2\times$ ($>4$ mm / $>1.0$ N)};
\node[mmproc] (g6) at (10.8,-7.2)   {6 Digital-twin discrepancy};
\node[mmdark] (aud) at (10.8,-2.4)  {Hash-chained audit trail\\$-24$h to $+72$h\\(21 CFR part 11)};
\draw[mmedge] (ev) to[out=60,in=180,looseness=0.6] (c7.west);
\draw[mmedge] (ev) -- (c15.west);
\draw[mmedge] (ev) to[out=-80,in=180,looseness=0.5] (g1.west);
\draw[mmedge] (ev) to[out=-90,in=180,looseness=0.5] (g2.west);
\draw[mmedge] (ev) to[out=-95,in=180,looseness=0.5] (g5.west);
\draw[mmedge] (g3) -- (aud.south);
\draw[mmedge] (g4) to[out=90,in=-90,looseness=0.7] (aud.south);
\draw[mmedge] (g6) to[out=90,in=-30,looseness=0.7] (aud.south);
\end{mermaidfig}
\figcaption{Figure 16. The safety-reporting clocks and the six Physical AI reporting
triggers. A fatal or life-threatening suspected adverse reaction is reported within
7 calendar days and any other serious and unexpected suspected adverse reaction
within 15 calendar days (21 CFR \S 312.32). The six Physical AI triggers of
\S 312.32(g) are a serious Physical AI adverse event (7- or 15-day), a system-drug
interaction (15-day), a cybersecurity incident (15-day, or 7-day if patient harm is
possible), sustained model degradation (at least 3 consecutive procedures or 24
cumulative hours), a sim-to-real divergence exceeding twice the validated tolerance
(trajectory over 4 mm or force over 1.0 N), and a digital-twin discrepancy. Each
preserves the hash-chained audit trail from minus 24 to plus 72 hours under 21 CFR
part 11.}

\subsubsection*{Discontinuation and stopping rules}

The protocol distinguishes discontinuation of an intervention from withdrawal from
the study; a participant who stops one or both components for a protocol-defined
reason remains in the study for safety and survival follow-up unless consent for that
follow-up is withdrawn. The investigational device is discontinued for the index
procedure, with conversion to the appropriate fallback technique, whenever any of
five prespecified conditions occurs: a conversion to open or manual technique judged
by the supervising surgeon to better serve patient safety or oncologic adequacy
(conversion is always available and is never itself a protocol violation); repeated
excursions against the $\leq 3$ N per-arm or $\leq 18$ N cumulative caps; any breach
of a vascular no-fly zone (1.0 mm at the superior mesenteric and portal veins;
1.5 mm at the hepatic artery, celiac axis, and superior mesenteric artery) or a
recurring pattern of no-fly approaches; any sim-to-real divergence beyond twice the
validated tolerance (trajectory deviation greater than 4 mm or force discrepancy
greater than 1.0 N), which also triggers a \S 312.32(g) report; and any hard-stop
cross-arm E-stop event, after which device resumption requires that the cause be
identified and the system re-pass the pre-procedure safety matrix. Daraxonrasib is
discontinued in an individual participant for a drug-attributable DLT during the
28-day window, for any other unacceptable toxicity not resolving with protocol dose
interruption and reduction, for a confirmed ISGPS Grade C fistula (which precludes
restart and maps to the T+21d or continued-hold arm of the advisory), for pregnancy,
for contraindicating intercurrent illness, or at the participant's request
\cite{onpremwhippl,daraxwhipple}.

Two binding halt rules govern the study as a whole and are applied by the data
safety monitoring board (DSMB): a device- or procedure-related serious adverse event
in at least one of three participants within a cohort, or two device-related deaths
at any time, triggers DSMB review and a halt. On any hold or closure, enrollment and
patient-facing robotic activity stop immediately, affected participants are
transitioned to appropriate clinical care, and a controlled decommissioning and
data-preservation procedure seals the complete hash-chained command and telemetry
record for the regulatory retention period. Table~\ref{tab:sec06-stopping}
consolidates the device discontinuation conditions, the drug discontinuation
conditions, and the two study-level halt rules, with the level at which each applies
and the action that follows, so that the entire stopping logic can be read at one
glance and audited against the source documents.

\needspace{10\baselineskip}
\begin{table}[t]
\caption{Table 6.11. Discontinuation and stopping rules, by the level at which each
applies. Device and drug discontinuations are individual-participant or
index-procedure rules; the two binding halt rules are study-level and rest with the
DSMB. Conversion to open or manual technique is always available and is never itself
a protocol violation.}
\label{tab:sec06-stopping}
\centering
\begin{tabularx}{\textwidth}{L{2.6cm} Y L{3.0cm}}
\toprule
\textbf{Level} & \textbf{Condition} & \textbf{Action} \\
\midrule
Device (index procedure) & Surgeon-judged conversion to open or manual technique for safety or oncologic adequacy & Convert; not a protocol violation \\
Device (index procedure) & Repeated excursions against the $\leq 3$ N per-arm or $\leq 18$ N cumulative force caps & Discontinue device; convert \\
Device (index procedure) & Breach of a vascular no-fly zone or a recurring pattern of no-fly approaches & Discontinue device; convert \\
Device (index procedure) & Sim-to-real divergence beyond $2\times$ tolerance (over 4 mm or over 1.0 N) & Discontinue device; \S 312.32(g) report \\
Device (index procedure) & Any hard-stop cross-arm E-stop event & Halt; resume only after cause identified and safety matrix re-passed \\
Drug (participant) & Drug-attributable DLT in the 28-day window; other unacceptable toxicity; Grade C fistula; pregnancy; intercurrent illness; participant request & Discontinue daraxonrasib; continue safety and survival follow-up \\
Study (DSMB) & Device- or procedure-related SAE in at least 1 of 3 within a cohort & DSMB review and halt \\
Study (DSMB) & Two device-related deaths at any time & DSMB review and halt \\
\bottomrule
\end{tabularx}
\end{table}

\subsubsection*{Statistical considerations and oversight}

The statistical approach is descriptive and estimation-based, non-confirmatory, and
follows ICH E9 and the TRIPOD+AI reporting standard for the model-assisted advisory
component. No formal power calculation is performed; exact two-sided 95\%
Clopper-Pearson confidence intervals are reported for the binomial endpoints, and a
two-sided alpha of 0.05 is used only in a supportive sense with no multiplicity
adjustment. Independent safety oversight operates on three tiers: the DSMB reviews
accumulating safety data at each cohort boundary and at any triggered review and
holds halt authority; the Independent Safety Monitor (ISM) is designated for each
procedure and provides real-time case-level oversight with the authority to call a
stop; and the Physical AI Safety Review Committee meets on a cadence not exceeding
90 days (and ad hoc on any triggering event) to review the autonomy-bearing behavior
of the system, the verification-before-generation gate-surface decisions, the USL
rating, and any software change, and to recommend USL reassessment or oversight
changes. Seven Physical AI record types (the LLM advisory record, the robot command
record, the sensor and telemetry record, the safety-limit-event record, the
verification-before-generation gate-surface record, the human-override and
intervention record, and the binding hash-chained audit trail) are retained for the
period required by 21 CFR \S 312.57, at least two years after a marketing
application is approved or, if none is filed, two years after the investigation is
discontinued and the FDA is notified \cite{daraxwhipple,clintrustpai}.

The benchmarking of the secondary endpoints against the 2025 Dutch nationwide
robotic pancreaticoduodenectomy cohort is descriptive and is not a formal comparison;
it situates the observed rates against a contemporary real-world reference so that an
unexpectedly high conversion, complication, fistula, or mortality rate is recognized
promptly even though the small sample cannot support a hypothesis test against the
reference. The seven Physical AI record types are summarized in
Table~\ref{tab:sec06-records} with their content and retention basis, so that the
records environment can be confirmed to satisfy 21 CFR \S 312.57 and 21 CFR part 11.
Each record type is append-only and hash-chained where applicable, so that the
provenance of every advisory verdict, every robot command, and every human override
is reconstructable for the full retention period, and so that the FDA can audit the
autonomy-bearing behavior of the system after the fact with the same fidelity as a
contemporaneous observer \cite{clintrustpai,paisitedocpk}.

\needspace{10\baselineskip}
\begin{table}[t]
\caption{Table 6.12. The seven Physical AI record types retained under 21 CFR
\S 312.57 and 21 CFR part 11, with the content and retention basis of each. Records
are retained for at least two years after a marketing application is approved or, if
none is filed, two years after the investigation is discontinued and the FDA is
notified.}
\label{tab:sec06-records}
\centering
\begin{tabularx}{\textwidth}{L{4.4cm} Y}
\toprule
\textbf{Record type} & \textbf{Content and retention basis} \\
\midrule
LLM advisory record & Each advisory verdict (ACCEPT, BLOCK, ESCALATE) with the hash-pinned seed and commit; retained for autonomy-bearing audit \\
Robot command record & Every issued command with the per-arm and cumulative force state at the heartbeat boundary \\
Sensor and telemetry record & The 640-channel stream (100 kHz force; 10 kHz other) including trajectory accuracy against the 0.05 mm specification \\
Safety-limit-event record & Every vascular-zone verdict, force-cap excursion, heartbeat or watchdog event, and E-stop activation \\
Verification-before-generation gate-surface record & The ten-gate suite decisions surrounding each advisory \\
Human-override and intervention record & Every manual override, conversion, and clinician confirmation \\
Binding hash-chained audit trail & The append-only chain that binds all of the above; preserved $-24$h to $+72$h around each Physical AI adverse event \\
\bottomrule
\end{tabularx}
\end{table}

The full protocol, version 1.0.0, with its complete schedule of activities,
statistical analysis plan, and Physical AI System Description meeting the eleven
fields of 21 CFR \S 312.23(g), is the bound supporting document for this section and
is archived at DOI 10.5281/zenodo.20780121
(\href{https://doi.org/10.5281/zenodo.20780121}{https://doi.org/10.5281/zenodo.20780121}).

\subsection{Informed Consent}

Informed consent is obtained by a qualified member of the research team who is not
the treating surgeon, before any study-specific procedure, in language the
participant can understand and with adequate time and opportunity to ask questions,
consistent with 21 CFR part 50 and ICH E6(R3) \cite{cfr050paiuni,ichpaistduni}. The
consent process incorporates each of the basic and additional elements of informed
consent required by 21 CFR \S 50.25, and a copy of the form and any participant
information sheet are submitted with the protocol. The consent discussion covers the
diagnosis, the alternatives (including standard open or conventional robotic
pancreaticoduodenectomy without the investigational advisory layer), and the
prognosis of a malignancy with a total five-year survival below 13\%; the
perioperative daraxonrasib regimen, its once-daily oral administration, the
protocol-mandated operative pause, and the toxicities characteristic of the RAS(ON)
inhibitor class; and, in addition to the standard elements, an explicit and distinct
disclosure of the on-premises LLM-directed eight-arm robotic system, its role as a
second-opinion advisory oracle held outside the vendor kinematic stack, its
continuous human oversight with immediate manual override, and its system-specific
risks. The voluntary nature of participation, the right to withdraw at any time
without penalty and without loss of any benefit to which the participant is otherwise
entitled, and the maintenance of confidentiality under the HIPAA Safe Harbor
de-identification method are stated at consent and reaffirmed at each contact.

The consent process is mapped explicitly to the basic and additional elements of 21
CFR \S 50.25 in Table~\ref{tab:sec06-consent}, so that an IRB reviewer can confirm
that each required element is addressed and can see which elements carry
Physical-AI-specific content beyond what a drug-only consent would contain. The
additional elements specific to this trial class are the distinct disclosure of the
autonomy-bearing device and its system-specific risks, the explanation of the
staged-autonomy model, and the documented Physical AI opt-out, each of which is
addressed in the dedicated paragraphs below. The consent form is written at an
accessible reading level, and the consent conversation is conducted with adequate
time, with qualified interpreters and translated materials where needed, so that the
consent is meaningful and not merely a signature \cite{cfr050paiuni,clintrustpai}.

\needspace{10\baselineskip}
\begin{table}[t]
\caption{Table 6.13. Mapping of the informed-consent process to the basic and
additional elements of 21 CFR \S 50.25, with the Physical-AI-specific content that
distinguishes this consent from a drug-only consent.}
\label{tab:sec06-consent}
\centering
\begin{tabularx}{\textwidth}{L{4.6cm} Y}
\toprule
\textbf{Consent element (21 CFR \S 50.25)} & \textbf{Content in this trial} \\
\midrule
Purpose and investigational nature & First-in-human dose-finding and early-feasibility study; primary purpose is safety and feasibility, not individual benefit \\
Reasonably foreseeable risks & Daraxonrasib toxicities of the RAS(ON) class; major pancreaticobiliary surgical risks; system-specific device risks \\
Alternatives & Standard open or conventional robotic pancreaticoduodenectomy without the investigational advisory layer \\
Confidentiality & HIPAA Safe Harbor de-identification; records available to the FDA \\
Voluntary participation and withdrawal & Right to withdraw at any time without penalty or loss of entitled benefit \\
Additional: device disclosure & Distinct disclosure of the on-premises LLM-directed eight-arm system as a second-opinion advisory oracle outside the vendor kinematic stack \\
Additional: staged autonomy & Clinician-controlled and supervised levels only at this phase; full autonomy prohibited per 21 CFR \S 312.21(e) \\
Additional: Physical AI opt-out & Documented right to the operative intervention without the Physical AI advisory layer where clinically feasible, without penalty \\
\bottomrule
\end{tabularx}
\end{table}

The consent form addresses the therapeutic misconception directly. Because the
population is vulnerable, facing the most technically demanding curative-intent
abdominal operation for a lethal cancer, the form and the consent conversation make
explicit that the study is a first-in-human, dose-finding, and early-feasibility
investigation whose primary purpose is to characterize safety and feasibility rather
than to confer individual therapeutic benefit, and that the LLM directs no autonomous
action. The staged-autonomy model is explained: in this Phase 1 the system operates
only at the clinician-controlled (teleoperation or supervised) and supervised
semiautonomous levels under continuous human oversight, full autonomy is prohibited
consistent with 21 CFR \S 312.21(e), and any future move to higher autonomy is
reserved for a later phase and gated by independent review. An independent consent
monitor is used when indicated, and qualified interpreters and translated materials
are provided so that consent is meaningful for participants with limited English
proficiency \cite{clintrustpai,oncotrialllm}.

A defining feature of this trial class is the documented Physical AI opt-out. Under
21 CFR \S 312.60(f), as adapted for Physical AI investigations, the participant is
offered, and may exercise, a documented right to request administration of the
operative intervention without the Physical AI advisory layer where clinically
feasible, without penalty and without affecting eligibility for the remainder of the
study or for clinical care \cite{cfr312paiuni}. The participant's election or
declination of the opt-out is recorded in the source document and the case report
form, and the schedule of activities (Table~\ref{tab:sec06-soa}) captures that
election at the consent visit. Where an unanticipated problem changes the risk
profile of the combined intervention, including a relevant Physical AI safety event,
the consent document is updated and currently enrolled participants are re-consented
before any further study procedure, consistent with the original opt-out.

\subsection{Investigator and Facilities Data}

The investigation is conducted at a single academic medical center through its
institutional multidisciplinary pancreatic-cancer program. Eligible individuals are
identified at tumor-board review and at surgical and medical-oncology clinics, and
the study is described to them by a member of the research team who is not the
treating surgeon, to reduce the potential for undue influence. The site possesses an
established high-volume pancreaticobiliary surgical service, a robotic surgical
program, an investigational-product pharmacy with controlled, temperature-monitored
storage and a drug-accountability log, an anatomic-pathology service applying a
standardized synoptic margin protocol with masked independent adjudication, a central
imaging capability for RECIST version 1.1 reads by a blinded reviewer, and the
clinical laboratory capacity for the serum daraxonrasib trough assay and the
hematologic, hepatic, renal, and pancreatic safety panels. The device facilities
include the qualified clinical environment in which the eight-arm platform is
maintained and calibrated between procedures, the operator control stations, the
on-premises compute environment that hosts the repository-pinned advisory LLM in
network isolation from the vendor kinematic stack under a deny-by-default
authentication and authorization model, and the 21 CFR part 11 records environment
that maintains the hash-chained, append-only audit trail. Day-to-day monitoring is
delegated to a robotics- and artificial-intelligence-competent clinical research
organization (CRO) over the site team, which comprises the operating surgeon, a
primary operator, and a qualified backup operator ready to assume control or invoke
the manual override at any time \cite{onpremwhippl,paisitedocpk}.

The facilities are listed against their function and their governing regulatory
basis in Table~\ref{tab:sec06-facilities}, so that the adequacy of the facilities
required by 21 CFR \S 312.23(a)(6)(iii)(b) can be confirmed for both the drug and the
device components. The site is selected for its combination of a high-volume
pancreaticobiliary surgical service and an established robotic program, because the
combined drug-device intervention demands surgeons and operators experienced in both
the open fallback technique and the robotic approach, and because the early-feasibility
nature of the device requires a site capable of maintaining, calibrating, and isolating
the on-premises compute environment to the standard required by 21 CFR part 11. The
network isolation of the advisory LLM from the vendor kinematic stack under a
deny-by-default model is a facility-level control that enforces the architectural
separation between the advisory oracle and the controller, so that the prohibition on
autonomous control is realized in the infrastructure and not only in the protocol
\cite{onpremwhippl,natlpaiplatf}.

\needspace{10\baselineskip}
\begin{table}[t]
\caption{Table 6.14. The clinical and device facilities supporting the investigation,
with the function and the governing regulatory basis of each. The facilities satisfy
the adequacy requirement of 21 CFR \S 312.23(a)(6)(iii)(b) for both the drug and the
device components.}
\label{tab:sec06-facilities}
\centering
\begin{tabularx}{\textwidth}{L{4.4cm} Y L{2.8cm}}
\toprule
\textbf{Facility} & \textbf{Function} & \textbf{Regulatory basis} \\
\midrule
High-volume pancreaticobiliary surgical service & Performs the pancreaticoduodenectomy and the open fallback & 21 CFR \S 312.23(a)(6) \\
Robotic surgical program & Hosts and operates the eight-arm platform & 21 CFR \S 812 (IDE) \\
Investigational-product pharmacy & Controlled, temperature-monitored storage; drug-accountability log & 21 CFR \S 312.57 \\
Anatomic-pathology service & Standardized synoptic margin protocol with masked adjudication & R0 endpoint \\
Central imaging capability & RECIST v1.1 reads by a blinded reviewer & Exploratory PFS/OS \\
Clinical laboratory & Serum daraxonrasib trough assay; safety panels & Safety assessments \\
On-premises compute environment & Hosts the repository-pinned advisory LLM in network isolation under deny-by-default & 21 CFR \S 312.21(e) \\
21 CFR part 11 records environment & Maintains the hash-chained, append-only audit trail & 21 CFR part 11 \\
\bottomrule
\end{tabularx}
\end{table}

The governance structure that supervises these investigators and facilities is the
three-tier oversight architecture described in \S 6.1: the Sponsor-Investigator
answers to a reviewing IRB and to the independent DSMB, convenes the Physical AI
Safety Review Committee on a cadence not exceeding 90 days, designates the
Independent Safety Monitor for each procedure, and delegates monitoring to the CRO.
The governance and safety-oversight structure that binds the investigators, the
facilities, and the independent reviewers together is presented in full in \S 10.4
(Figure~20), and is summarized for the present section in the oversight table below.

The three-tier oversight architecture is summarized in
Table~\ref{tab:sec06-oversight}, with the membership, the review cadence, and the
authority of each tier. The tiers are deliberately complementary in their temporal
scope: the Independent Safety Monitor provides real-time, case-level oversight during
each procedure with immediate stop authority; the DSMB provides periodic,
cohort-level oversight at each cohort boundary and at any triggered review with
binding halt authority; and the Physical AI Safety Review Committee provides
autonomy-focused oversight on a cadence not exceeding 90 days and ad hoc on any
triggering event, with authority to recommend USL reassessment and oversight changes.
The combination ensures that an emerging safety signal is caught at whichever
timescale it manifests, from a single intraoperative event to a slow degradation
across procedures, and that the binding stop and halt authorities rest with bodies
independent of the conduct of the study \cite{daraxwhipple,clintrustpai}.

\needspace{9\baselineskip}
\begin{table}[t]
\caption{Table 6.15. The three-tier independent safety-oversight architecture, with
the membership, review cadence, and authority of each tier. The binding stop and halt
authorities rest with bodies independent of the conduct of the study.}
\label{tab:sec06-oversight}
\centering
\begin{tabularx}{\textwidth}{L{3.6cm} Y L{3.0cm}}
\toprule
\textbf{Tier} & \textbf{Membership and cadence} & \textbf{Authority} \\
\midrule
Independent Safety Monitor (ISM) & Designated for each procedure; independent of the operating team; real-time & Stop authority at the case level \\
Data and Safety Monitoring Board (DSMB) & Surgical, oncologic, biostatistical, and device-safety members; at each cohort boundary and any triggered review & Binding halt authority; applies the two halt rules \\
Physical AI Safety Review Committee & Reviews autonomy-bearing behavior, gate-surface decisions, USL, and software changes; cadence not exceeding 90 days and ad hoc & Recommends USL reassessment and oversight changes \\
\bottomrule
\end{tabularx}
\end{table}

\subsubsection{FDA Form 1572}

A Statement of Investigator (FDA Form 1572) is executed by each investigator
responsible for the conduct of the study at the clinical site, in satisfaction of
21 CFR \S 312.53(c), and is assembled with the submission together with the
investigator's current curriculum vitae and the documentation of relevant
qualifications. The executed Form 1572 attests that the investigator has the
training and experience to assume responsibility for the proper conduct of the
investigation, identifies the name and address of the clinical site and of the
clinical laboratory and other facilities to be used, names the reviewing IRB
responsible for review and approval of the study, lists the sub-investigators (the
operating surgeon, the primary and backup device operators, the medical oncologist,
the pathologist, and the research coordinators) who will assist in the conduct of
the investigation, and identifies the protocol by its title and version. By signing,
the investigator commits to conduct the study in accordance with the protocol, to
personally conduct or supervise the investigation, to inform participants that the
drug and device are being used for investigational purposes and to ensure that
consent and IRB requirements are met, to report adverse events to the sponsor in
accordance with 21 CFR \S 312.64, to read and understand the investigator's brochure
and the Physical AI System Description, to ensure that all associates and staff are
informed of their obligations, to maintain adequate and accurate records available
for inspection under 21 CFR \S 312.68, and to comply with all other requirements of
21 CFR part 312. The executed forms, with attachments, are bound into the submission
and are updated whenever a sub-investigator or facility changes \cite{phase1guidance}.

The commitments made by signing Form 1572 are summarized against their regulatory
basis in Table~\ref{tab:sec06-1572}, so that the completeness of the investigator's
undertaking can be confirmed at a glance and so that the additional commitment
specific to this trial class, to read and understand the Physical AI System
Description as well as the investigator's brochure, is made visible alongside the
conventional commitments. The Form 1572 commitments are the mechanism by which the
sponsor's obligations are devolved to the investigator at the site, and the executed
forms together with the curricula vitae establish that the investigation will be
conducted by qualified persons as required by 21 CFR \S 312.23(a)(6)(iii)(b)
\cite{phase1guidance,paisitedocpk}.

\needspace{9\baselineskip}
\begin{table}[t]
\caption{Table 6.16. The investigator commitments made by executing FDA Form 1572,
with the regulatory basis of each. The commitment to read and understand the
Physical AI System Description is specific to this trial class.}
\label{tab:sec06-1572}
\centering
\begin{tabularx}{\textwidth}{Y L{3.4cm}}
\toprule
\textbf{Commitment} & \textbf{Regulatory basis} \\
\midrule
Conduct the study in accordance with the protocol and personally conduct or supervise the investigation & 21 CFR \S 312.60 \\
Inform participants that the drug and device are investigational; ensure consent and IRB requirements are met & 21 CFR \S\S 312.60, 50.25 \\
Report adverse events to the sponsor & 21 CFR \S 312.64 \\
Read and understand the investigator's brochure and the Physical AI System Description & 21 CFR \S 312.23(g) \\
Ensure all associates and staff are informed of their obligations & 21 CFR \S 312.53(c) \\
Maintain adequate and accurate records available for inspection & 21 CFR \S 312.68 \\
Comply with all other requirements of 21 CFR part 312 & 21 CFR part 312 \\
\bottomrule
\end{tabularx}
\end{table}

\subsubsection{Disclosure of Financial Interests}

Financial interests are disclosed and managed in accordance with 21 CFR part 54,
Financial Disclosure by Clinical Investigators, and the certification and disclosure
statements (FDA Forms 3454 and 3455 as applicable) are collected before participation
and updated through one year after completion of the study. The disclosure makes
explicit the principal conflict of interest in this investigation: the
Sponsor-Investigator, Kevin Kawchak, is the chief executive officer of
ChemicalQDevice, the entity that develops the on-premises LLM-directed eight-arm
robotic Physical AI system under study and holds the combined IND and IDE. This
relationship is disclosed in full to the FDA and to the reviewing IRB, and it is
managed by the independent oversight structure described in \S 6.3 so that safety and
escalation decisions are not made by a party with a financial interest in the outcome
alone \cite{cfr312paiuni,clintrustpai}.

The mitigation is structural and operates at every decision point that could be
influenced by the conflict. The DSMB, composed of members independent of the conduct
of the study with surgical, oncologic, biostatistical, and device-safety expertise,
holds the binding halt authority and applies the prespecified halt rules (a device-
or procedure-related serious adverse event in at least one of three participants
within a cohort, or two device-related deaths at any time). The Independent Safety
Monitor, designated for each procedure and independent of the operating team, holds
real-time stop authority at the case level. The Physical AI Safety Review Committee
reviews the autonomy-bearing behavior of the system, the verification-before-generation
gate-surface decisions, the USL rating, and any software change on a cadence not
exceeding 90 days, and recommends USL reassessment or oversight changes. Routine
monitoring is delegated to a CRO independent of the device manufacturer, and audits
may be conducted by a designee independent of the conduct of the study. Because the
binding safety, escalation, and halt decisions rest with these independent bodies
rather than with the Sponsor-Investigator, the financial conflict is disclosed and
mitigated to a degree commensurate with the first-in-human, autonomy-bearing nature
of the investigation, and any identified conflict is disclosed and managed before it
can affect the conduct, analysis, or reporting of the study \cite{daraxwhipple}.

The conflict and its structural mitigation are summarized in
Table~\ref{tab:sec06-coi}, with the independent body to which each binding decision
is assigned, so that an FDA or IRB reviewer can confirm that no binding safety,
escalation, or halt decision rests with the conflicted Sponsor-Investigator. The
guiding principle is that disclosure alone is insufficient for an autonomy-bearing,
first-in-human investigation in which the Sponsor-Investigator develops and
commercializes the device under study; the conflict is therefore managed by
relocating every binding decision to an independent body, and by delegating routine
monitoring to a CRO independent of the device manufacturer. The financial-disclosure
certifications are updated through one year after completion of the study under 21
CFR part 54, so that any change in the financial relationships is captured for the
full reporting period \cite{cfr312paiuni,paiautospons}.

\needspace{9\baselineskip}
\begin{table}[t]
\caption{Table 6.17. The principal financial conflict of interest and its structural
mitigation, with the independent body to which each binding decision is assigned. No
binding safety, escalation, or halt decision rests with the conflicted
Sponsor-Investigator.}
\label{tab:sec06-coi}
\centering
\begin{tabularx}{\textwidth}{L{4.4cm} Y}
\toprule
\textbf{Binding decision} & \textbf{Independent body and basis} \\
\midrule
Study-level halt & DSMB, independent of the conduct of the study; applies the two prespecified halt rules \\
Case-level stop & Independent Safety Monitor, independent of the operating team; per-procedure real-time authority \\
Autonomy and USL changes & Physical AI Safety Review Committee; cadence not exceeding 90 days and ad hoc \\
Routine monitoring & CRO independent of the device manufacturer \\
Audit & Designee independent of the conduct of the study \\
Financial disclosure & 21 CFR part 54; Forms 3454/3455; updated through one year after completion \\
\bottomrule
\end{tabularx}
\end{table}
